\chapter{Marco teórico}

Los algoritmos de ordenamiento y búsqueda son fundamentales en el procesamiento de datos, ya que permiten organizar información y acceder a ella de manera eficiente. Su estudio no solo considera la forma en que funcionan, sino también el costo computacional asociado a su ejecución, el cual suele analizarse en función del tamaño de la entrada.

En este trabajo se utilizan algoritmos de ordenamiento iterativos clásicos: Bubble Sort, Insertion Sort, Selection Sort y Cocktail Shaker Sort. Estos métodos permiten reordenar los registros de deportistas según distintos criterios, como ID, nombre, equipo, puntaje o cantidad de competencias. Aunque todos cumplen la misma función general, presentan diferencias en su rendimiento dependiendo de si la entrada se encuentra ordenada, desordenada o invertida.

Por otra parte, para la localización de registros se emplean búsqueda secuencial y búsqueda binaria iterativa. La búsqueda secuencial revisa los elementos uno a uno hasta encontrar el valor buscado o llegar al final del arreglo, mientras que la búsqueda binaria divide repetidamente el espacio de búsqueda, por lo que requiere que los datos estén previamente ordenados.

El análisis de complejidad permite estimar el comportamiento de estos algoritmos en mejor, peor y caso promedio. Esta comparación teórica sirve como base para interpretar los resultados experimentales obtenidos durante la ejecución de los benchmarks realizados en el sistema.
